\begin{proof}[Proof of \cref{obj:lem:rho-deficient-regime-algebra}]
\emph{Step 1 (strict comparison of the two exponents).} Since \(s>0\) we have \(4s>0\), so the
deficient-regime hypothesis \(4s<d\) can be divided through:
\[
1<\frac d{4s}.
\]
Adding \(1+1/\alpha\) to both sides,
\[
2+\frac1\alpha
<
1+\frac d{4s}+\frac1\alpha .
\]
Both sides are strictly positive, because \(1/\alpha>0\) by \(\alpha>0\). Dividing the constant
\(2>0\) by the two sides therefore reverses the strict inequality:
\[
\frac{2}{1+d/(4s)+1/\alpha}
<
\frac{2}{2+1/\alpha}.
\]
Multiplying numerator and denominator by \(\alpha>0\) gives the elementary identity
\[
\frac{2}{2+1/\alpha}
=
\frac{2\alpha}{2\alpha+1},
\]
and hence the strict exponent comparison asserted in the statement,
\[
\frac{2}{1+d/(4s)+1/\alpha}
<
\frac{2\alpha}{2\alpha+1}.
\]
% lean: rho_deficient_regime_algebra

\emph{Step 2 (identification of the maximum in \(\rho_n\)).} If \(n=1\), every real power of \(1\)
equals \(1\), so both terms of \cref{obj:def:published-hoif-rate} coincide,
\[
1^{-2\alpha/(2\alpha+1)}=1=1^{-2/(1+d/(4s)+1/\alpha)},
\]
and \(\rho_1=\max\{1,1\}=1=1^{-2/(1+d/(4s)+1/\alpha)}\), as claimed. If instead \(n>1\), the map
\(x\mapsto n^{x}\) is strictly increasing; negating the two exponents reverses the strict inequality
of Step 1, and applying \(x\mapsto n^{x}\) then preserves the reversed inequality, so
\[
n^{-2\alpha/(2\alpha+1)}
\le
n^{-2/(1+d/(4s)+1/\alpha)} .
\]
In either case the maximum in \cref{obj:def:published-hoif-rate} is attained by the
covariate-smoothness term, that is,
\[
\rho_n
=
n^{-2/(1+d/(4s)+1/\alpha)} .
\]
Together with the strict comparison of Step 1, this is the desired conclusion.
% lean: rho_deficient_regime_algebra
\end{proof}
